\section{Introduction} 

Write $\mathbb N = \{1,2,\ldots\}$ and $\mathbb N_0 = \{0, 1,2,\ldots\}$. 
A subset $A$ of the integers is said to have \emph{positive upper Banach density} if 
\begin{equation}  \label{e:banach}
    d ^{\ast}(A) \coloneqq \limsup_{N-M\to\infty}  \frac{\lvert A\cap [M,N) \rvert} {N-M} >0. 
\end{equation} 
A set $B \subset \mathbb N_0$ is \emph{syndetic} if it has bounded gaps, that is, if there is an $L$ so that 
$B \cap [n, n+L] \neq \emptyset $ for every $n$.  A syndetic set has positive lower density, at least $(L+1) ^{-1}$; the converse is false.  

Recall that an \emph{IP-set} is built out of an infinite sequence of  
 integers $r_1  ,  r_2 ,\ldots $,  which need not be distinct: 
the IP-set is the collection of all the finite sums $ \sum_{n\in F} r_n $, 
where $F\subset \mathbb N$ is finite and non-empty.  
The theorem below is the IP Szemer\'edi theorem of Furstenberg and Katznelson 
\cite{FK85}*{Theorem A}, specialized to the powers of a single transformation.  

\begin{priorResults}  \label{t:FK}
    Let $(X, \mathcal A, m, T)$ be an invertible measure preserving system, let $A \in \mathcal A$ have $m(A)>0$, and let $\ell \geq 1$ be an integer. 
    Then, for every IP-set $\mathcal I$, there is an $r \in \mathcal I$ so that 
    \begin{equation}
        m\bigl(A \cap T ^{-r} A \cap \cdots \cap T ^{-\ell r} A \bigr) >0. 
    \end{equation}
\end{priorResults}

By the Furstenberg correspondence principle, Theorem \ref{t:FK} is already a \emph{density} statement about the integers.  
Namely, if $A \subset \mathbb Z$ has positive upper Banach density, and $\mathcal I$ is an IP-set, 
then there is an $r \in \mathcal I$ for which the set of $a$ with 
$a, a+r , \ldots, a+ \ell r \in A$ again has positive upper Banach density.  
Indeed, it follows that the set of \emph{good differences}
\begin{equation} \label{e:good}
\mathsf{G} = \mathsf{G}(A , \ell ) \coloneqq 
\bigl\{ r \in \mathbb N \colon d ^{\ast} \bigl( A \cap (A-r) \cap \cdots \cap (A - \ell r) \bigr) >0 \bigr\} 
\end{equation}
meets every IP-set, that is, $\mathsf G$ is an IP$^{\ast}$ set.  
Such a set is syndetic: the complement of a set which is not syndetic contains arbitrarily long intervals, and a set which contains arbitrarily long intervals contains an IP-set.  

\smallskip 

The sets of integers of interest to us are these.  
Let $\mathcal{C}_{b,D}$ denote a \emph{classical integer Cantor set}, defined by a base $b\geq 3$ 
and a collection of digits 
$ D \subset \{0,1 , \ldots, b-1\}$ with $ 2\leq \lvert D \rvert < b$.  
That is, 
\begin{equation} \label{e:Cantor}
\mathcal{C}_{b, D}  \coloneqq 
\Bigl\{ 
 \sum_{j=0}^k d_j b ^{j} 
 \colon  k\in \mathbb{N}_0 ,\ d_j \in  D
\Bigr\}. 
\end{equation}
The most classical version is the middle thirds Cantor set $\mathcal{C}_{3,\{ 0, 2 \}}$, that is base $b=3$ 
and allowed digits $D= \{0,2\}$.  
Throughout the paper we assume that $0 \in D$, and we write 
$\mathcal C = \mathcal C_{b,D} = \{ 0=k_0 < k_1 < k_2 < \cdots\}$ for the increasing enumeration of the Cantor set.  

A Cantor set with $0 \in D$ contains an IP-set.  Indeed, fixing a non-zero $d \in D$, 
the finite sums of the integers $d b ^{j}$, $j \geq 0$, all have their base $b$ digits in $\{0,d\} \subset D$.  
Therefore, Theorem \ref{t:FK} applies to the IP-set inside $\mathcal C$, and so already provides 
elements of $\mathcal C$ which are good differences.  
The same remark applies to the recent work of Kra and Shalom \cite{kra2025ergodicaverageslargeintersection}, 
who show that for IP-sets with rational spectra---a class which includes the IP-set just described---the  
IP averages above have the same characteristic factors as the full linear averages, and who obtain 
the optimal lower bounds for progressions of length three and four.  
It should be emphasized that the implication runs from IP-sets to Cantor sets, and not the other way:  
the Cantor set is the \emph{larger} set, so that locating a good difference inside the IP-set that it contains is the stronger statement.  

\smallskip 

What these results do not address is the \emph{frequency} with which the good differences occur inside the Cantor set.  
A Cantor set comes with its own natural enumeration $\{k_n \colon n \geq 0\}$, and the question we take up is how often 
$k_n$ is a good difference, as measured in the index $n$.  

It should be said at once that the two properties of $\mathsf G$ recorded above do not answer this question.  
Syndeticity of $\mathsf G$ carries no information about $\mathsf G \cap \mathcal C$, since $\mathcal C$ has density zero: 
a syndetic set of integers can be disjoint from a Cantor set, as one sees by taking $s = \gcd (D)>1$, so that $\mathcal C \subset s \mathbb Z$, 
and taking the syndetic set of integers not divisible by $s$.  
Nor does the IP$^{\ast}$ property settle it.  That property places an element of $\mathsf G$ inside every \emph{infinite} IP-set.  
What is wanted here is, for some fixed $i$, an element of $\mathsf G$ in each of the blocks 
$b ^{i}k_u + \bigl( \mathcal C \cap [0,b ^{i}) \bigr)$, $u \geq 0$, into which Lemma \ref{l:Diff}(ii) decomposes $\mathcal C$.  
These blocks are finite, and they are translated.  

\smallskip 

Our main theorem is the following.  It is stated in terms of averages over intervals of indices 
$M \leq n < N$, with $N - M \to \infty$; the uniformity in $M$ is exactly what delivers the frequency statement.  

\begin{theorem}  \label{t:CSZ}
    Let $\mathcal C = \mathcal C_{b,D}=\{0=k_0<k_1 < \cdots\}$ be a Cantor set with  $0\in D$.  
    Then, for all integers $\ell \geq 1 $, all  invertible  measure preserving systems $(X,\mathcal A, m, T)$, 
    and  all non-negative  $f \in L ^{\infty}(X)$ which are not identically zero, we have 
    \begin{equation} \label{e:CSZ}
        \liminf_{N-M\to\infty} \Expectation_{M \leq n < N} 
        \int_X \prod_{j=0}^{\ell}  T ^{j k_n} f  \;dm >0. 
    \end{equation}
    \end{theorem}

\begin{corollary} \label{c:syndetic}
    With the notation and assumptions of Theorem \ref{t:CSZ}, the set of indices 
    \begin{equation}
    \Bigl\{ n \in \mathbb N_0 \colon \int_X \prod_{j=0}^{\ell}  T ^{j k_n} f  \;dm >0 \Bigr\} 
    \end{equation}
    is syndetic.  
\end{corollary}

\begin{proof}
    Write $c>0$ for the left hand side of \eqref{e:CSZ}, and suppose that the set of indices above is not syndetic.  
    Then there are intervals $[M_i, N_i)$ with $N_i - M_i \to \infty$ which contain no such index.  
    The average in \eqref{e:CSZ} over $[M_i,N_i)$ is then zero for every $i$, contradicting \eqref{e:CSZ}.  
\end{proof}

Passing to the integers, we obtain the frequency statement that is the point of this paper.  

\begin{theorem}  \label{t:CSZ-Integers}
    Let $\mathcal C =\{0=k_0<k_1 < \cdots\}$ be a Cantor set with $0\in D$, and let $\ell \geq 1$ be an integer.  
    For every $A \subset \mathbb Z$ of positive upper Banach density, the set of indices 
    \begin{equation} \label{e:B}
        B \coloneqq \{ n \in \mathbb N \colon k_n \in \mathsf G (A, \ell ) \} 
        = \bigl\{ n \in \mathbb N \colon d ^{\ast} \bigl( A \cap (A-k_n) \cap \cdots \cap (A - \ell k_n) \bigr) >0 \bigr\} 
    \end{equation}
    is syndetic.  In particular, for every $n \in B$, the set $A$ contains a non-degenerate $(\ell +1)$-term arithmetic progression with step size $k_n$.  
    \end{theorem}

\begin{proof}
By the Furstenberg correspondence principle \cite{MR0498471}*{Thm 1.1}, there is an invertible measure preserving system 
$(X,\mathcal A, m ,T)$ and a set $\tilde A \in \mathcal A$ with $m (\tilde A) = d ^{\ast}(A)>0$, such that for every finite collection of integers $r_1 , \ldots, r_j$, 
\begin{equation} \label{e:correspondence}
d ^{\ast} \bigl( A \cap (A - r_1) \cap \cdots \cap (A-r_j) \bigr) 
\geq 
m \bigl( \tilde A \cap T ^{-r_1} \tilde A \cap \cdots \cap T ^{-r_j} \tilde A \bigr) . 
\end{equation}
We apply \eqref{e:correspondence}  for each fixed index $n$, with $(r_1 , \ldots, r_ \ell ) = (k_n , 2k_n , \ldots, \ell k_n)$.  
It follows that $B$ contains the set of indices $n \geq 1$ for which 
$m ( \tilde A \cap T ^{-k_n} \tilde A \cap \cdots \cap T ^{- \ell k_n} \tilde A )>0$, and the latter is syndetic by 
Theorem \ref{t:CSZ} and Corollary \ref{c:syndetic}, applied to $f = \mathbf 1_{\tilde A}$.  
(Discarding the index $n=0$, for which $k_0=0$ and the statement is trivial, does not affect syndeticity.)  
\end{proof}

\begin{remark}  \label{r:order}
    The correspondence principle is applied above for one index $n$ at a time, and this is essential.  
    The averages over the index $n$ are taken \emph{after} the passage from the integers to the measure preserving system, 
    and it is only the syndeticity of the set of good indices which survives that passage.  
    An average over $x$ and $n$ taken simultaneously, as in 
    $\liminf_{N} \Expectation_{x \in [N]} \Expectation_{n \in [N]} \prod_{j=0}^{\ell}\mathbf 1_{A} (x - jk_n)$, 
    is not a formal consequence of Theorem \ref{t:CSZ}: 
    the limit in $x$ is taken first in the correspondence principle, and in general 
    $\lim_{n} \lim _{m} a_{m,n} = a$ does not imply $\lim_n a_{n,n}=a$.  
\end{remark}

\begin{remark} \label{r:novelty}
    It is useful to separate the several questions that arise here.  
    That a set of positive density contains $(\ell+1)$-term progressions is Szemer\'edi's theorem.  
    That the step size may be required to lie in a prescribed IP-set, and that a positive density of such progressions is then obtained, 
    is the theorem of Furstenberg and Katznelson, Theorem \ref{t:FK}; by the remarks above it covers Cantor sets as well.  
    The optimal lower bounds for the density of progressions, with step size restricted to a rational IP-set, are due to Kra and Shalom, 
    for progressions of length three and four.  
    The question addressed here is instead one of frequency: among the step sizes $k_n$ that a Cantor set makes available, 
    how many are good differences?  Theorem \ref{t:CSZ-Integers} answers that a syndetic set of them are.  
    Frequency questions of this kind seem to have received relatively little attention.  
\end{remark}


Our proof of Theorem \ref{t:CSZ} follows the maximal-factor strategy of the ergodic proof of Szemer\'{e}di's Theorem \cites{MR0498471,MR670131}, 
carried out so that all averages are uniform in the index, as in \eqref{e:CSZ}.  
Section \ref{s:prelim} sets this up.  The CSZ property \eqref{e:CSZdef} of a measure preserving system carries a dilation parameter, which the compact case requires, 
and Theorem \ref{t:CSZ} follows from three statements: the property passes to weakly mixing and to compact extensions of a CSZ factor (Lemma \ref{l:Extensions}); 
it is preserved by suprema of chains of factors (Lemma \ref{l:chain}); and every proper factor has an extension of one of these two kinds (Lemma \ref{l:dichotomy}).  
These give the result for ergodic systems on compact metric spaces, and Proposition \ref{p:reduction} reduces the general case to that one.  
The arithmetic of the Cantor set enters through two facts: the block structure of its enumeration (Lemma \ref{l:Diff}), 
and a uniform estimate for exponential sums $\sum e(\beta k_n)$ over intervals of indices (Lemma \ref{l:irrational}).  

The Cantor set imposes a real constraint at two steps.  For weakly mixing extensions, in \S\ref{sec:wm}, the van der Corput argument leads to correlations 
along the differences $k_{n+h}-k_n$; these take few values, so the correlation for an individual shift $h$ is not small.  
We bound each shift by a non-negative quantity, and only then average over the shifts, using the spectral measure of $f\otimes \overline f$ on the 
relatively independent self-joining.  For compact extensions, in \S\ref{sec:compact}, the argument of Furstenberg that passes through van der Waerden's theorem 
produces return times which are dilates, and a Cantor set is not closed under dilation.  We follow instead the argument of Furstenberg, Katznelson and Ornstein \cite{MR670131}, 
whose return times are translates, which a Cantor set does absorb.  A Cantor set version of van der Waerden's theorem is proved in \S\ref{sec:vdW}, independently of the rest.  

Further questions that this result presents to us include: 
\begin{itemize}
    \item The question of quantitative bounds for these sorts of configurations is certainly challenging, and could be of some interest. 
  In the special case of $\ell=1$,  examples of power savings bounds are contained in \cite{firstpaper}. 
  
    \item  Again, for $\ell=1$, there is a polynomial version of this result \cite{firstpaper}. The result above should surely hold for polynomial configurations as well. 
    
    \item  The higher dimensional analog could have different formulations as well.  
    
\end{itemize}

\begin{ack}
    We thank Or Shalom for a careful reading of an earlier version of this paper, and for the comments which led to the
    present formulation of Theorems \ref{t:CSZ} and \ref{t:CSZ-Integers}.  
\end{ack}
